\section{Discussion}
\label{sec:discussion}

\subsection*{Identity and classification under equivalence}
The organising principle of modern and contemporary mathematics is
classification up to structural sameness: objects are known not in
themselves but through the isomorphisms that relate them.  This has a long
prehistory in the rhetorical structure of Vitruvius's \emph{De
Architectura}, whose opening dichotomy \emph{quod significatur et quod
significat}---the signified and the signifier, the \emph{res} and the
\emph{verba}---makes the work knowable through the proportional relations
its \emph{modulus} mediates: by number, proportion, and anatomy, in the
pre-Cartesian mode of interpretation Foucault
describes~\cite{santos-vitruvio,foucault-1966}.  The modulus is one unit
from which the proportions of the whole are read---identity generating
equivalence.

Modern mathematics inherits precisely this tension.  Two objects may be
\emph{equal}, literally the same, or merely \emph{isomorphic}, the same up
to a structure-preserving relabelling; classification proceeds almost always
by the second.  The present note is unusual in sitting at the meeting of the
two.  Its centre is a genuine identity---a single $\golden$, not a class of
isomorphic objects, satisfies twelve clauses at once---yet the clauses are
organised by equivalence: the residue classes of the Galois group
$\Ztwenty$, the split and inert types in $\mathbb{Q}(\sqrt5)$.  The same
$\golden$ is read through the equivalences that classify the structures it
inhabits.  The arc below moves between the two registers---identity where
the equation holds by construction, equivalence where the classification
does the work---answering the millennial question of
\Cref{ssec:classification} in the specific case of~$\golden$.

The note can be read as a single arc, from a tautology to a spectral
endpoint, each step naming where it is established.

\subsection*{The simultaneous setting}
The frame of the note is that a single $\golden$ is inscribed at once in
five canonical structures---the complex rotor of Euler's identity, the real
trigonometric line, the real multiplicative line, the arithmetic of Mersenne
modulo~$20$, and the Galois group of $\mathbb{Q}(\zeta_{20})/\mathbb{Q}$%
---the twelve simultaneous clauses of \Cref{thm:coincidence-intro}.  This is
the golden-ratio counterpart of Euler's identity: where $e^{i\pi}+1=0$ knots
analysis, algebra, and geometry at one point, here the same $\golden$ knots
those structures through the pentagon $\golden=2\cos(\pi/5)$ and its
cyclotomic inscription $\golden=\zeta_{10}+\zeta_{10}^{-1}$.  The central
identity below is $\golden$'s inscription on the multiplicative line; the
arc that follows traces its depth.

\subsection*{The tautological identity}
Read as an equality, the central identity $3\,\golden^{\sigma(p)}=2^{p}$ is
tautological: the golden index
$\sigma(p):=p\,\lambda_{\log}-\log_{\golden}3$ (\Cref{def:sigma}) is defined
precisely so that it holds, for every real~$p$, and its proof
(\Cref{thm:universal}) is the evaluation of that definition.  Nothing about
the truth of the equation is at stake; it holds by construction.  Yet the
slope it carries, $\lambda_{\log}=\log_{\golden}2$, is transcendental over
$\mathbb{Q}$ (Gelfond--Schneider, \Cref{rem:transcendence})---the exact
counterpart of the period $2\pi$ in Euler's identity, an algebraic relation
between algebraic numbers carried by a transcendental exponent.

\subsection*{Its non-tautological content}
What is not tautological is what the constituents of the equation turn out
to be.  The multiplier~$3$ is not a free constant: it equals
$M_{2}=|\{3,7,11\}|$, the number of residue classes the Mersenne primes
occupy modulo~$20$ (\Cref{thm:concentration-main},
\Cref{thm:three-main})---a fact about the arithmetic of $2^{p}-1$ that owes
nothing to the equation.  Here the identity first carries content beyond
its definition.

\subsection*{Its proof from outside}
That the three constituents---the factor~$3$, the base~$\golden$, the slope
$\lambda_{\log}$---are each pinned by a theorem \emph{external} to the
identity, and are mutually consistent with the golden coordinate uniquely
determined, is made precise in \Cref{thm:characterization}: the identity is
a consistency theorem for independently fixed determinations, certified
without circularity, not a discovered equality.  This is the sense in which
it is proved from outside itself.

\subsection*{Its expansion into the Mersenne primes}
Read forward, into the primes it produces, the identity reaches the 52
Mersenne exponents.  On the multiplicative line their binary powers
$2^{p}=\Mp+1$ fit a single integer sublattice of the golden base, every pair
standing in an exact golden ratio (\Cref{thm:ratio-main})---in the binary
guise, the shift of the single carried bit from one position to another
(\Cref{rem:binary})---with 52 distinct transcendental indices $\sigma(p_{i})$
(\Cref{cor:linear-independence}).

In the field $\mathbb{Q}(\sqrt5)$ the reach is finer, and it resolves each
Mersenne prime into presentations of one invariant.  The exponent residue
$p\bmod4$ fixes at once its class in $\Ztwenty$, its splitting type
(\Cref{cor:split-from-p}), and its Fibonacci sign $F_{\Mp}\equiv\pm1$
(\Cref{thm:fib-main}):
\[
  p\equiv1\mapsto(11,\ \mathrm{split},\ +1),\quad
  p\equiv3\mapsto(7,\ \mathrm{inert},\ -1),\quad
  p=2\mapsto(3,\ \mathrm{inert},\ -1),
\]
the triple recording the class modulo~$20$, the type in $\mathbb{Q}(\sqrt5)$,
and the sign of $F_{\Mp}$; and the inert primes are exactly those with
$\Mp\mid F_{\Mp+1}$ (\Cref{prop:rank-inert}).  Over the 52 known primes the
joint presentation takes exactly three values, one per admissible residue
(\Cref{cor:joint}); given $p\bmod4$, no further information separates the
class, the type, and the sign.  The development is \Cref{sec:fibonacci}.

\subsection*{The spectral meeting point}
Read backward, the same period $\lambda_{\log}$ carries $\golden$ down to
one half, $\golden^{-\lambda_{\log}}=\tfrac12$ (\Cref{sec:spectral}).  This
is the value at which the companion operator paper anchors its
Hilbert--P\'olya operator---the point at which the golden arithmetic
developed here meets the spectral programme.

\medskip
Throughout, the note works on the axis of representation and equivalence
rather than literal equality.  What a comparison of expressions registers as
distinct---$\Mp$ against $2^{p}$, an index defined against one proved---the
relating equivalences return to a single object in several guises.  This is
quantitative, not figurative: given the exponent residue $p\bmod4$, the
residue class, the splitting type, and the Fibonacci sign carry no
information beyond one another, so that over the 52 known Mersenne primes
their joint presentation takes exactly three values, one per admissible
residue.  The content lies not in any single presentation but in the exact
maps between them; where a literal reading sees distinct objects that fail to
coincide, the equivalence measures no difference at all.

